\section{Introduction}

Machine learning progress is increasingly measured through benchmark performance. Papers With Code (PWC) hosts thousands of benchmark leaderboards where research teams compete to achieve state-of-the-art (SOTA) results on shared evaluation datasets. Yet benchmarks themselves have lifecycles: a benchmark that defines the frontier today may be superseded tomorrow when the community adopts a harder or more comprehensive successor. Understanding what predicts this \emph{benchmark displacement} --- the event where one benchmark loses plurality status to another within the same task --- has implications for how we design evaluation infrastructure, allocate research effort, and interpret progress claims.

A natural candidate predictor is \emph{community breadth at benchmark introduction}: how many distinct research teams participated in the benchmark when it first achieved plurality status. Two competing mechanisms motivate this hypothesis. Under the \textbf{lock-in mechanism} (H1), broad early adoption creates a large network of stakeholders who have invested in the benchmark --- trained models, established baselines, published results. This investment creates switching costs, slowing displacement (hazard ratio HR $<$ 1). Under the \textbf{saturation mechanism} (H2), widespread adoption signals that the community has collectively ``solved'' the benchmark, accelerating the search for harder successors and replacement (HR $>$ 1). Both mechanisms are theoretically grounded: H1 follows Rogers' Diffusion of Innovations framework for opinion-leader lock-in \cite{rogers2003diffusion}; H2 aligns with the ICLR 2025 workshop framing of benchmark overuse and community-driven replacement \cite{iclr2025benchmarking}.

The challenge in testing this hypothesis is measurement. Prior attempts using raw cumulative submission counts collapsed into time proxies (partial $r^{2}$ = 0.0011 with temporal controls) --- a methodological failure that renders Cox regression uninterpretable. A prior directional signal (HR = 0.871 for \emph{score velocity}, a different construct) appeared only with 22 displacement events --- an underpowered sample where sampling noise dominates. The current paper tests \emph{submission-count diversity} at adequate power; the two results are complementary, not contradictory (see Section~\ref{sec:prior}).

This paper makes three contributions:

\begin{enumerate}
  \item \textbf{Methodological:} We introduce a 5-gate FAIL FAST pre-validation protocol (G0--G4) that separates data quality from hypothesis validity before any regression. Gates test join coverage (G0), time-independence of diversity predictors via partial $R^{2}$ (G1--G2), cross-benchmark variance (G3), and multicollinearity (G4). All five gates pass by wide margins, ruling out measurement failure as an explanation for the null.

  \item \textbf{Empirical:} We conduct the most statistically powered test to date of community breadth \emph{as operationalized by paper submission count} as a displacement predictor. On 258 complete-case rows from the h-e2 panel (87 Koch et al.\ \cite{koch2021reduced} taxonomy tasks, 2015--2023, Papers With Code), with events-per-variable (EPV) $\approx$ 86, we find HR = 1.006 (95\% CI = [0.846, 1.196], LRT $p$ = 0.9495). This is a \textbf{near-perfect null}.

  \item \textbf{Scientific:} The combination of validated measurement and adequate power means the null is scientifically informative. Community breadth diversity --- operationalized as unique paper submission count at introduction year --- is ruled out as a predictor class for benchmark displacement timing in the Papers With Code ecosystem (2015--2023). Score velocity and institutional diversity remain untested at full power.
\end{enumerate}
